\section{Metatheory} %
\label{sec:meta}

In this section, we provide the denotational semantics of types
($\denot{\tau}$) and type contexts ($\denot{\Gamma}$), and state the
soundness theorem of our type system.  The key challenge is to give a
semantics that supports both demonic and angelic modalities
simultaneously, together with the resource-based interpretation of
contexts that tracks variable dependencies.  We proceed in four
steps: we define contextual capabilities as the semantic domain
(Section~\ref{subsec:capabilities}), equip them with algebraic structure
(Section~\ref{subsec:algebra}), build a context logic over this algebra
(Section~\ref{subsec:logic}), and finally give the type denotation and
soundness theorem (Section~\ref{subsec:denotation}).

\subsection{Contextual Capabilities}%
\label{subsec:capabilities}

In a pure functional language, the free variables of a term $e$ are
bound by its evaluation environment, a partial map from variables to
values.  To reason about both safety and reachability, we need to
consider not just a single environment but a \emph{set} of
environments, representing the range of possible variable bindings
available to demonic or angelic choices.  We call such a set a
\emph{contextual capability} or capability for short.

\begin{definition}[Contextual Capability]\label{def:contextual-capability}
  A \emph{contextual capability} is a non-empty set of environments
  (partial maps, $\Var \rightharpoonup \Val$) that all share the same
  domain.  We use $r$ to range over capabilities and $\Res$ as the set
  of all capabilities.
\end{definition}

The domain condition ensures that all environments in a capability
bind exactly the same set of variables, reflecting the fact that the
free variables of a term are fixed by its typing context.  A
capability with a single environment corresponds to a fully determined
variable binding; a capability with multiple environments encodes the
range of choices available for safety or reachability reasoning.

We extend the standard notions of environment compatibility
($r_1 \uparrow r_2$, meaning every environment in $r_1$ is compatible
with every environment in $r_2$), projection ($\projectTo{r}{X}$, the
piecewise restriction of $r$ to the variables in a set $X$), and
refinement ($r_1 \sqsubseteq r_2$, meaning
$r_1 = \projectTo{r_2}{\Dom{r_1}}$) from environments to capabilities
in the obvious pointwise way.

The refinement relation $\sqsubseteq$ plays a central role: $r_1
\sqsubseteq r_2$ says that $r_1$ is a restriction of $r_2$ to a
smaller variable domain, i.e., $r_2$ extends the bindings of $r_1$
with additional variables.  This relation will serve as the Kripke
accessibility relation in our logic (Section~\ref{subsec:logic}).

To express variable dependency in refinement type qualifiers, we also
need the following notion.

\begin{definition}[Fiber] %
  \label{def:fibers}
  For a variable $x$ and an environment $[x \mapsto v] \in
  \projectTo{r}{\{x\}}$, the \emph{fiber} of $r$ over $[x \mapsto
  v]$ is:
  \[
    \fiberOver{r}{[x \mapsto v]} \doteq \{\, \sigma \in r \mid
    \sigma(x) = v \,\}.
  \]
  That is, $\fiberOver{r}{[x \mapsto v]}$ is the largest subset of $r$
  comprised of environments that map $x$ to $v$.
\end{definition}
\noindent Fibers allow us to condition a capability on the value of a
specific variable, which is essential for evaluating refinement
qualifiers that refer to other variables bound in the typing context
(see Section~\ref{subsec:logic}).

\subsection{Capability Algebra} %
\label{subsec:algebra}

Our context logic is a bunched implication (BI) logic~\cite{OHP99}
based on a BI-algebra~\cite{Pym+02}, a partial commutative monoid with
a pre-order satisfying bifunctoriality.  We equip contextual
capabilities with two operations: a \emph{product} that composes
disjoint variable bindings, and a \emph{sum} that merges capabilities
over the same domain to support control-flow reasoning. The
\emph{product} $r_1 \times r_2$ operation, defined when
$r_1 \uparrow r_2$, is the set of all pairwise unions
$\sigma_1 \cup \sigma_2$ with $\sigma_1 \in r_1$ and
$\sigma_2 \in r_2$.
%% it combines capabilities over disjoint variable domains. \ZZ{the varaible domain is allowed to be not disjointed, which is important for resource duplication which will be used in sec 4 and 5. }
The \emph{sum} operation, $r_1 + r_2$, defined when
$\Dom{r_1} = \Dom{r_2}$, is simply $r_1 \cup r_2$; it merges
capabilities over the same domain to support control-flow
reasoning. Together with unit $\mathbf{1} \doteq \{\emptyset\}$, the
tuple $(\Res, +, \times, \mathbf{1})$ forms a partial commutative
semiring without zero. To enable a Kripke semantics for BI-logic, this
semiring must be equipped with a partial order satisfying
bifunctoriality.

The natural candidates for this order are the subset ($\subseteq$) and
superset ($\supseteq$) relations.  However, neither works for both
angelic and demonic modalities simultaneously: the subset order forces
an overapproximation semantics (demonic), while the superset order
forces underapproximation (angelic).  Since we need to support both as
first-class operators, the Kripke order must be \emph{neutral} with
respect to approximation.  The refinement relation $\sqsubseteq$ has
exactly this property: $r_1 \sqsubseteq r_2$ holds if $r_2$ extends
$r_1$ with additional bindings, while preserving all existing
environments, neither adding nor removing possible values for the
variables already in $r_1$'s domain.  Specifically, $\sqsubseteq$
satisfies the standard bifunctoriality property required of a
BI-algebra: $r_1 \sqsubseteq r_1'$ and $r_2 \sqsubseteq r_2'$ implies
$r_1 \times r_2 \sqsubseteq r_1' \times r_2'$, ensuring that
$(\Res, +, \times, \mathbf{1}, \sqsubseteq)$ forms a BI-algebra.


\subsection{Context Logic} %
\label{subsec:logic}

We now define a logic over our capability algebra.  The logic is a BI
logic extended with demonic and angelic modalities, a sum operator for
control-flow reasoning, and a \emph{binding reference operator} that
expresses how a qualifier may depend on specific variable bindings in
the current capability.

\begin{definition}[Context Logic]\label{def:context-logic}
  The syntax of our context logic is:
  \begin{align*}
    P, Q \;::=\; &\chtrue \mid \chfalse \mid P \chland Q \mid P \chlor Q \mid P \chimpl Q \mid P \chstar Q \mid P \chwand Q
    \\& \mid \chforall x.\, P \mid \chexists x.\, P \mid \Code{Atom}(\phi)
    \mid \demonic P \mid \angelic P \mid P \choplus Q \mid \forallM{x} P
  \end{align*}
  The forcing semantics over capabilities $r \in \Res$ is:
  \begin{itemize}
  \item $r \models \chtrue$ always; \quad $r \models \chfalse$ never;
  \item $r \models P \chland Q$ iff $r \models P$ and $r \models Q$;
  \item $r \models P \chlor Q$ iff $r \models P$ or $r \models Q$;
  \item $r \models P \chimpl Q$ iff for all $r'$ such that $r \sqsubseteq r'$, $r' \models P$ implies $r' \models Q$;
  \item $r \models P \chstar Q$ iff $\exists r_1, r_2.\; r_1 \times r_2 \sqsubseteq r$,\; $r_1 \models P$,\; $r_2 \models Q$;
  \item $r \models P \chwand Q$ iff $\forall r' \uparrow r$,\; $r' \models P$ implies $r' \times r \models Q$;
  \item $r \models \chforall x.\, P$ iff $x \notin \Dom{r}$ and $\forall r'$ such that $r \sqsubseteq r'$ with $x \in \Dom{r'}$,\; $r' \models P$;
  \item $r \models \chexists x.\, P$ iff $x \notin \Dom{r}$ and $\exists r'$ such that $r \sqsubseteq r'$ with $x \in \Dom{r'}$,\; $r' \models P$;
  \item $r \models \Code{Atom}(\phi)$ iff $\projectTo{r}{\Code{FV}(\phi)} = \{\,\sigma \mid \Dom{\sigma} = \Code{FV}(\phi) \land \sigma \models \phi\,\}$;
  \item $r \models \demonic P$ iff $\exists r', r \subseteq r'$ such that $r' \models P$;
  \item $r \models \angelic P$ iff $\exists r', r' \subseteq r$ such that $r' \models P$;
  \item $r \models P \choplus Q$ iff $\exists r_1, r_2.\; r_1 + r_2 \sqsubseteq r$,\; $r_1 \models P$,\; $r_2 \models Q$;
  \item $r \models \forallM{x} P$ iff $\forall [x \mapsto v] \in \projectTo{r}{\{x\}}.\; \fiberOver{r}{[x \mapsto v]} \models P[x \mapsto v]$.\footnote{
We extend the binding operator to sets of variables by iterated
application: for a finite set $X = \{x_1, \ldots, x_n\}$,
we write $\forallM{X}\, P$ as shorthand
for
$\forallM{x_1}\,\forallM{x_2}\,\cdots\,\forallM{x_n}\,P$
In particular,
$\forallM{\Dom{\Sigma}}\, P$ binds all
variables in the domain of the current basic type context $\Sigma$
before evaluating $P$.}
  \end{itemize}
  We write $P \models Q$ iff for all $r \in \Res$, $r \models P$ implies $r \models Q$.
\end{definition}

The additive connectives ($\chland$, $\chlor$, $\chimpl$, $\chforall$,
$\chexists$) are standard.  The multiplicative connectives ($\chstar$,
$\chwand$) follow standard BI-logic semantics using our algebra's
product.  We discuss the remaining operators below.

\paragraph{Kripke Monotonicity}
The implication $P \chimpl Q$ and the first-order quantifiers
$\chforall x.\, P$ and $\chexists x.\, P$ use the refinement relation
$\sqsubseteq$ as the Kripke accessibility relation: truth at $r$
requires truth at all $r'$ such that $r \sqsubseteq r'$, i.e., at all
extensions of $r$ with additional variable bindings.  This is the
standard Kripke semantics for implication in BI-logic, and ensures the
following monotonicity property.

\begin{theorem}[Kripke Monotonicity]\label{thm:kripke-monotonicity}
  $r \sqsubseteq r'$ and $r \models P$ implies $r' \models P$.
\end{theorem}

The use of $\sqsubseteq$ rather than $\subseteq$ or $\supseteq$ as the
accessibility relation is what allows our logic to remain neutral with
respect to approximation direction, per \autoref{subsec:algebra}.

\paragraph{Atomic Propositions}
The atomic proposition $\Code{Atom}(\phi)$ lifts a standard formula
$\phi$ to a capability-level predicate with \emph{exact} semantics: a
capability $r$ satisfies $\Code{Atom}(\phi)$ iff its projection to
$\Code{FV}(\phi)$ contains \emph{exactly} those environments that
satisfy $\phi$.  This neutrality leaves the choice of approximation
direction entirely to the demonic and angelic modalities.

\paragraph{Demonic and Angelic Modalities}
The demonic modality $\demonic P$ holds at $r$ if there exists an
upper bound $r'$ (i.e., $r \subseteq r'$) that satisfies $P$; this
models an overapproximation since $r'$ extends $r$ with additional
environments.  Dually, the angelic modality $\angelic P$ holds at $r$
if some lower bound $r' \subseteq r$ satisfies $P$; this models an
underapproximation since $r'$ witnesses a specific subset of the
environments in $r$.
\begin{example}%[Modalities as Lower- and Upper-Bounds]%
  \label{ex-approximation-modalities}
  Let $r_1$ be $\{[x \mapsto 1], [x \mapsto 2]\}$. The following hold:
  \begin{align*}
    &r_1 \models \demonic\,\Code{Atom}(x = 1) \;\judgfail \quad
    &&r_1 \models \demonic\,\Code{Atom}(0 \leq x \leq 4) \;\judgok \\
    &r_1 \models \angelic\,\Code{Atom}(x = 1) \;\judgok \quad
    &&r_1 \models \angelic\,\Code{Atom}(0 \leq x \leq 4) \;\judgfail
  \end{align*}
  The demonic assertion $\demonic\,\Code{Atom}(x = 1)$ fails because
  no superset of $r_x$ satisfies $\Code{Atom}(x = 1)$ exactly (since
  $r_x$ already contains $[x \mapsto 2]$, any superset violates
  exactness).  The angelic assertion
  $\angelic\,\Code{Atom}(0 \leq x \leq 4)$ fails because no subset of
  $r_x$ satisfies $\Code{Atom}(0 \leq x \leq 4)$ exactly (that would
  require environments for $x = 0, 3, 4$ which $r_x$ does not
  contain).

  Critically, satisfiability of demonic and angelic modalities is also
  maintained by the refinement relation $\sqsubseteq$. To illustrate,
  consider the following set of capabilities
\begin{align*}
\begin{array}{c @{\hspace{1.5em}} c @{\hspace{1.5em}} c}
  {\boxed{r_1 \doteq \{[x \mapsto 1], [x \mapsto 2]\}}}
  & \sqsubseteq
  & {\boxed{r_2 \doteq \{[x \mapsto 1; y \mapsto 2], [x \mapsto 2; y \mapsto 7]\}}} \\[0.6em]
  \vsubseteq & & \vsubseteq \\[0.6em]
  {\boxed{r_3 \doteq \{[x \mapsto 1]\}}}
  & \sqsubseteq
  & {\boxed{r_4 \doteq \{[x \mapsto 1; y \mapsto 2]\}}}
\end{array}
\end{align*}\noindent
For any angelic formula $\angelic P$ where $r_3 \models \angelic P$,
then it is also the case $r_2 \models \angelic P$. We have a similar
result for the demonic formula $\demonic P$, as
$r_1 \models \demonic P$ implies $r_4 \models \demonic P$.  Under the
later modality, the subset relation is reversed; however, the
refinement relation is not.

Finally, to illustrate how the product operation interacts with the
logic's two conjunction operators, let
$r_y \doteq \{[y \mapsto 3], [y \mapsto 4]\}$ and
$r_{xy} \doteq \{[x \mapsto 1, y \mapsto 3], [x \mapsto 2, y \mapsto
4]\}$.  Then:
\begin{align*}
&r_x \times r_y \models (\angelic\,\Code{Atom}(1 \leq x \leq 2)) \chstar (\angelic\,\Code{Atom}(3 \leq y \leq 4)) \;\judgok \\
&r_{xy} \models (\angelic\,\Code{Atom}(1 \leq x \leq 2)) \chland (\angelic\,\Code{Atom}(3 \leq y \leq 4)) \;\judgok \\
&r_{xy} \models (\angelic\,\Code{Atom}(1 \leq x \leq 2)) \chstar (\angelic\,\Code{Atom}(3 \leq y \leq 4)) \;\judgfail
\end{align*}
The product $r_x \times r_y$ satisfies the multiplicative conjunction
because $x$ and $y$ are independently bound.  The capability $r_{xy}$
satisfies the additive conjunction but not the multiplicative one,
since its bindings are entangled and cannot be decomposed into two
disjoint sub-capabilities.

\end{example}

\paragraph{Sum Operator}
The sum $P \choplus Q$ holds at $r$ if $r$ can be split into two
sub-capabilities $r_1 + r_2 \sqsubseteq r$, one satisfying $P$ and one
satisfying $Q$.  This supports the context sum type
$\tau_1 \ctoplus \tau_2$ from in \autoref{sec:context-typing}, which
allows a capability to be partioned into two subsets. The requirement
that capabilities be non-empty ensures that neither branch is vacuous:
both summand must contain at least one environment.

\paragraph{Binding Reference Operator}
The binding reference operator $\forallM{x} P$ expresses how a
qualifier may depend on the specific binding of variable $x$ in the
current capability.  It holds at $r$ if, for every possible binding
$v$ of $x$ in $r$, the sub-capability $\fiberOver{r}{[x \mapsto v]}$
satisfies $P$ with $x$ substituted by $v$.  This is needed because
$\Code{Atom}(\phi)$ has exact semantics: a qualifier $\phi$ mentioning
$x$ requires a satisfying capability $r$ to contain \emph{all}
environments consistent with $x$'s full range, which conflicts with
the constrained range of $x$ expressed by other conjuncts.  The
binding reference operator resolves this by fixing the binding of $x$
before evaluating the qualifier, enabling dependent qualifiers to
refer to specific variable bindings in the capability.

\begin{example}% [Binding Reference Operation]
  \label{ex-variable-reference}
  Consider the entangled context
  \[x{:}\urt{\Nat}{1 \leq \vnu
    \leq 2} \ctcomma y{:}\urt{\Nat}{0 \leq \vnu < x}\]
\noindent where $y$'s
  range depends on $x$.  Its denotation is:
  \[
    \denot{x{:}\urt{\Nat}{1 \leq \vnu \leq 2} \ctcomma y{:}\urt{\Nat}{0 \leq \vnu < x}}
    = \angelic\,(1 \leq x \leq 2) \chland \forallM{x}\,\angelic\,(0 \leq y < x).
  \]
  The operation $\forallM{x}\,\angelic\,(0 \leq y < x)$ fixes the
  binding of $x$ before evaluating $\angelic\,(0 \leq y < x)$, so
  $y$'s required range is determined separately for each value of $x$.
  Capabilities that model this formula include
  $\{[x \mapsto 1;\, y \mapsto 0]\}$ and
  $\{[x \mapsto 2;\, y \mapsto 0];\, [x \mapsto 2;\, y \mapsto 1]\}$,
  as well as their sum. Omitting $\forallM{x}$ from the right conjunct
  would require a satisfying capability to include an environment with
  a mapping for \emph{every} pair of $x$ and $y$ where $0 \leq y < x$:
  $[x \mapsto 3;\, y \mapsto 0], [x \mapsto 4;\, y \mapsto 1],
  \ldots$.
\end{example}

\subsection{Type Denotation and Soundness}%
\label{subsec:denotation}

We now give the denotation of types and type contexts in terms of the
context logic defined above, and state the main soundness theorem.  We
write $\denot{\tau}_\Sigma : \Code{Term} \to \Code{Formula}$ for the
type denotation under a basic type context $\Sigma$, and
$r \models \denot{\tau}_\Sigma\;e$ to mean that $e$ inhabits type
$\tau$ under capability $r$.  We omit $\Sigma$ when it is clear from
context.

As is customary in other refinement types, we require terms to be
total and well-typed in the basic type system.  We encode these side
conditions into the logic via a guard formula:
\begin{align*}
  P_{\Code{guard}} \doteq \forallM{\Dom{\Sigma}}\,\Code{Atom}(\Code{Total}(e) \land \Sigma \vdash_s e : \eraserf{\tau} \land \Sigma \vdash_{\Code{WF}} \tau).
\end{align*}
The binding reference operator $\forallM{\Dom{\Sigma}}$ is necessary
here; without it, e.g., $\Code{Atom}(\Sigma \vdash_s x : \Nat)$ is
only satisfied by a capability that contains an environment for every
possible value of $x$.  All type denotations below implicitly include
this guard, which we omit for brevity.

The denotation of a term $e$ under capability $r$ captures the image
of $e$ under all environments in $r$ via a variable $x$:
\begin{align*}
  \denot{\mstep{e}{x}}_\Sigma &\doteq \forallM{\Dom{\Sigma}}\,\Code{Atom}(\mstep{e}{x}).
\end{align*}
The binding reference operator $\forallM{\Dom{\Sigma}}$ fixes the
bindings of all variables in the typing context before evaluating the
atomic reduction predicate, ensuring that the reduction $\mstep{e}{x}$
is checked under each specific environment in the current capability
rather than requiring the capability to contain all environments
consistent with $x$'s full range.

\begin{definition}[Type Denotation]\label{def:type-denotation}
  \begin{align*}
    \denot{\ort{b}{\phi}}_\Sigma\;e &\doteq \denot{\mstep{e}{\vnu}}_\Sigma \chimpl \forallM{(\Code{FV}(\phi)-\{\vnu\})}\,\demonic\,\Code{Atom}(\phi) \quad \text{($\vnu$ fresh)}
    \\\denot{\urt{b}{\phi}}_\Sigma\;e &\doteq \denot{\mstep{e}{\vnu}}_\Sigma \chimpl \forallM{(\Code{FV}(\phi)-\{\vnu\})}\,\angelic\,\Code{Atom}(\phi) \quad \text{($\vnu$ fresh)}
    \\\denot{x{:}\tau_x \ctsarr \tau}_\Sigma\;e &\doteq \denot{\tau_x}_\Sigma\;x \chimpl \denot{\tau}_{\Sigma,x{:}\eraserf{\tau_x}}\;(e\;x) \quad \text{($x$ fresh)}
    \\\denot{x{:}\tau_x \ctwand \tau}_\Sigma\;e &\doteq \denot{\tau_x}_\Sigma\;x \chwand \denot{\tau}_{\Sigma,x{:}\eraserf{\tau_x}}\;(e\;x) \quad \text{($x$ fresh)}
    \\\denot{\tau_1 \ctinter \tau_2}_\Sigma\;e &\doteq \denot{\tau_1}_\Sigma\;e \chland \denot{\tau_2}_\Sigma\;e
    \\\denot{\tau_1 \ctunion \tau_2}_\Sigma\;e &\doteq \denot{\tau_1}_\Sigma\;e \chlor \denot{\tau_2}_\Sigma\;e
    \\\denot{\tau_1 \ctoplus \tau_2}_\Sigma\;e &\doteq \denot{\tau_1}_\Sigma\;e \choplus \denot{\tau_2}_\Sigma\;e
  \end{align*}
\end{definition}

Demonic and angelic type denotations follow directly from the
corresponding modalities: the result of evaluating $e$ bound to
fresh variable $\vnu$, must satisfy $\phi$ under either the demonic or
angelic interpretation, with the binding operator fixing the bindings
of the free variables of $\phi$ (besides $\vnu$) before applying the modality.  The
function type denotations internalize the adjunction between function
types and context combinators: $\ctsarr$ uses additive implication
(the parameter binding is entangled with the closure context), while
$\ctwand$ uses multiplicative implication (the parameter binding is
disjoint from the closure context).

\begin{definition}[Type Context Denotation]\label{def:context-denotation}
  \[
  \begin{aligned}
    \denot{\emptyset}_\Sigma &\doteq \chtrue \\
    \denot{x{:}\tau}_\Sigma &\doteq \denot{\tau}_\Sigma\;x
  \end{aligned}
  \qquad
  \begin{aligned}
    \denot{\Gamma_1 \ctcomma \Gamma_2}_\Sigma &\doteq \denot{\Gamma_1}_\Sigma \chland \denot{\Gamma_2}_{(\Sigma,\,\eraserf{\Gamma_1})} \\
    \denot{\Gamma_1 \ctstar \Gamma_2}_\Sigma &\doteq \denot{\Gamma_1}_\Sigma \chstar \denot{\Gamma_2}_\Sigma \\
    \denot{\Gamma_1 \ctoplus \Gamma_2}_\Sigma &\doteq \denot{\Gamma_1}_\Sigma \choplus \denot{\Gamma_2}_\Sigma
  \end{aligned}
  \]
\end{definition}

The dependent conjunction ``$\ctcomma$'' maps to additive conjunction
$\chland$, reflecting that the second context may refer to bindings in
the first.  The disjoint conjunction ``$\ctstar$'' maps to multiplicative
conjunction $\chstar$, requiring the two sub-capabilities to be
decomposable as a product — i.e., their variable bindings are
independent.

The semantic subtyping relations used in rules T-Sub and T-CtxSub are:
\begin{figure}[h]
  {
  \begin{flalign*}
    &\text{\textbf{Semantic subtyping}} && \fbox{$\Gamma \vdash \tau \subty \tau$ \quad $\Gamma \ctsubseteq{X} \Gamma$}
  \end{flalign*}
  \mprooftr{42mm}
  {$\forall e.\; \denot{\Gamma} \models \denot{\tau_1}\;e \chimpl \denot{\tau_2}\;e$}
  {Sub}
  {$\Gamma \vdash \tau_1 <: \tau_2$}
  \quad
  \mprooftr{66mm}
  {$\forall r.\; r \models \denot{\Gamma_1} \implies \exists r'.\; \projectTo{r}{X} = \projectTo{r'}{X} \land r' \models \denot{\Gamma_2}$}
  {CtxSub}
  {$\Gamma_1 \ctsubseteq{X} \Gamma_2$}
  }
  \caption{Semantic subtyping relations \BD{Can we add subscripts here
      or explicitly state the $\Sigma$ we're using in this caption.}}
  \label{fig:semantic-subtyping-relation}
\end{figure}

Type subtyping is encoded directly as logical implication in the
denotational semantics.  Context subtyping uses an explicit projection
$\projectTo{r}{X}$: the two capabilities agree on the variable set $X$
(the free variables of the term and type), but $r'$ may differ on
other variables, reflecting the fact that context subtyping only needs
to hold for the variables that actually appear in the judgement.

We conclude with the core soundness results.

\begin{definition}[Well-Formed Operator Context] %
  \label{thm:well-formed-operator}
  Operator context $\Phi$ is well-formed iff for all primitive
  operators $\primop$:
  {\small\begin{align*}
    \Phi(\primop) = \overline{x{:}\ort{b_x}{\phi_x}} \ctsarr \ort{b}{\phi} \ctinter \urt{b}{\phi}
    \quad\text{and}\quad
    \models \denot{\overline{x{:}\ort{b_x}{\phi_x}}} \chiff \denot{\ort{b}{\phi} \ctinter \urt{b}{\phi}}\;(\primop\;\overline{x}).
  \end{align*}}
\end{definition}

\BD{We do not mention the above definition anywhere. Presumably it's a
  premise of the theorems below?}

\begin{theorem}[Fundamental Theorem]%
  \label{thm:fundamental}
  The typing judgement $\Gamma \vdash e : \tau$ is sound with respect
  to the denotation of types in context logic:
  $\Gamma \vdash e : \tau$ implies
  $\denot{\Gamma}_\emptyset \models
  \denot{\tau}_{\eraserf{\Gamma}}\;e$. %Thus, any capability that
  %satisfies the denotation of $\Gamma$ also satisfies the denotation
  %of $\tau$ applied to $e$
\end{theorem}

\begin{corollary}[Type Soundness]%
  \label{thm:soundness}
  $\emptyset \vdash e : \tau$ implies that for any fresh variable $x$,
  $\{[x \mapsto v] \mid \mstep{e}{v}\} \models \denot{\tau}\;x$.
\end{corollary}

Our type soundness theorem simply instantiates the fundamental theorem
at the empty context. For a closed well-typed term $e$, the capability
$\{[x \mapsto v] \mid \mstep{e}{v}\}$, which collects all values $e$
can reduce to, each bound to a fresh variable $x$, satisfies the
denotation of $\tau$. This has two concrete consequences depending on
the modality of $\tau$: if $\tau$ is a demonic type, then every value
that $e$ reduces to satisfies $\tau$'s qualifier (safety); if $\tau$
is an angelic type, then every value prescribed by $\tau$'s qualifier
is witnessed by some reduction of $e$ (reachability).
